% ---
% school: "college"
% level: "8eme"
% section: null
% chapter: "العمليات والحساب على الأعداد الكسرية"
% source: "Axiomaths"
% title: "الحنفيات"
% country: "Tunisie"
% year: 2026
% professor: "Axiomaths"
% difficulty: "Moyen"
% points: null
% tags: ["fractions", "robinets", "probleme", "debit"]
% ---
% Exercice 11/18 de la série 8ème « العمليات والحساب على الأعداد الكسرية » (statement + \axsol + solution). Compiler avec XeLaTeX + axiomath.sty (+ axiomaths-logo.png).
\documentclass[12pt,a4paper]{article}
\usepackage[level=8b,color=15803D]{axiomath}
\axsetup{title={العمليات والحساب على الأعداد الكسرية}, corrige}

\begin{document}

\begin{corr}[title={الحنفيات}, stars=3, source={Axiomaths}]
خزان فارغ. الحنفية الأولى تملؤه وحدها في $6$ ساعات والحنفية الثانية تملؤه وحدها في $4$ ساعات.
\begin{enumerate}
  \item ما هو الكسر من الخزان الذي تملؤه كل حنفية في ساعة واحدة ؟
  \item نفتح الحنفيتين معا. ما هو الكسر من الخزان الذي يمتلئ في ساعة واحدة ؟
  \item استنتج المدة اللازمة لملء الخزان بالحنفيتين معا، بالساعات ثم بالساعات والدقائق.
  \item في أسفل الخزان ثقب يفرغه وحده في $12$ ساعة. نفتح الحنفيتين والثقب مفتوح. ما هي المدة اللازمة لملء الخزان ؟
\end{enumerate}
\axsol
\begin{enumerate}[itemsep=14pt]
  \item الحنفية الأولى تملأ الخزان في $6$ ساعات، فهي تملأ في ساعة واحدة $\dfrac{1}{6}$ منه.
  الحنفية الثانية تملأ في ساعة واحدة $\dfrac{1}{4}$ منه.
  \item الكسر الذي يمتلئ في ساعة بالحنفيتين معا :
  \[\dfrac{1}{6}+\dfrac{1}{4}=\dfrac{2}{12}+\dfrac{3}{12}=\dfrac{5}{12}\]
  \item في كل ساعة يمتلئ $\dfrac{5}{12}$ من الخزان، فلملء الخزان كله نحتاج :
  \[1\div\dfrac{5}{12}=1\times\dfrac{12}{5}=\dfrac{12}{5}\]
  أي $\dfrac{12}{5}$ ساعة، وهذا يساوي :
  \[\dfrac{12}{5}=2+\dfrac{2}{5}\]
  $\dfrac{2}{5}$ من الساعة تساوي $\dfrac{2}{5}\times60=24$ دقيقة. إذن المدة هي \ib{$2$ ساعة و $24$ دقيقة}.
  \item الثقب يفرغ في ساعة $\dfrac{1}{12}$ من الخزان. الكسر الذي يمتلئ في ساعة :
  \[\dfrac{5}{12}-\dfrac{1}{12}=\dfrac{4}{12}=\dfrac{1}{3}\]
  المدة اللازمة :
  \[1\div\dfrac{1}{3}=1\times3=3\]
  أي $3$ ساعات.
\end{enumerate}
\textbf{الفكرة :} نجمع (أو نطرح) ما ينجزه كل عنصر \emph{في وحدة زمن واحدة} ثم نأخذ المقلوب.
\end{corr}

\end{document}
